\chapter{Supervised Instruction Tuning}
\label{ch:supervised-instruction-tuning}

\abstract{This chapter explains how pretrained models become instruction-following models through supervised examples. It covers chat templates, task diversity, Alpaca-style self-instruction, synthetic data, refusal data, data quality, and the limits of imitation.}

\section{Chapter Spine}
Instruction tuning changes the interface a model expects. The chapter separates the historical importance of Alpaca-style recipes from the modern requirements of dataset auditing, chat formatting, safety coverage, and evaluation.

\section{Core Sections}
\subsection{Instruction Data}
Prompt-response pairs, multi-turn dialogue, system messages, tool traces, and formatting contracts.

\subsection{Synthetic Data}
Self-Instruct, distillation from stronger models, diversity filters, and error amplification.

\subsection{Training and Evaluation}
Loss masking, sequence packing, held-out tasks, refusal behavior, and benchmark leakage.

\subsection{Limits of SFT}
Why imitation alone does not produce calibrated, harmless, or preference-aligned behavior \cite{ouyang2022training}.

\section{Planned Exercises}
Convert a raw instruction dataset into a chat template and inspect which tokens should contribute to loss.
